\section{Results and Analysis}
\label{sec:results}

\noindent\emph{The tables in this section are structural placeholders. Every
numeric cell is marked \TBD{} and is populated only by the confirmatory run of
Section~\ref{sec:setup}; the development-validation runs are used solely to verify
that the analysis layer emits correct, well-formed metrics, and their numbers are
deliberately not reported here.}

The results below are computed across five valid conditions: which vocabulary,
\textit{free}, \textit{closed}, or \textit{open}, is paired with which
reasoning strategy, \textit{zero}, \textit{few}, or \textit{scientific}.
Figure~\ref{fig:conditions} (Section~\ref{sec:strategies}) lays out the full
space and marks the four invalid combinations. \textit{scientific-open} is the
default used for RQ1, RQ3, and RQ5; RQ2 and RQ4 additionally draw on the other
valid combinations, as noted in each subsection below. We settled on
\textit{scientific-open} as the default for two reasons. First, forcing every
bug into a fixed seven-type list manufactures false certainty for bugs that
genuinely do not fit, so the open vocabulary, with its mandatory justification
for every escape, turns taxonomy coverage into something we measure rather than
assume. Second, in early testing, simply narrating a step-by-step reasoning
style left the model's answers unchanged; only when that reasoning was
enforced, forcing a committed prediction before new evidence was revealed, did
a meaningful share of labels actually change, always moving toward the
fix-informed reference.

\subsection{RQ1: Defect-type Distribution}

RQ1 characterises the ODC composition of the benchmark from the pre-fix
\textit{scientific-open} classifications. Table~\ref{tab:rq1} reports the
distribution over the seven Design/Code types. The second half of RQ1, whether
the distribution varies significantly across projects, is tested with a
$\chi^2$ association once per-project counts are large enough to satisfy its
minimum expected-count assumption; at smaller per-project samples it is reported
descriptively.

\begin{table}[H]
\centering
\caption{RQ1: distribution of ODC defect types (pre-fix, \textit{scientific-open}).}
\label{tab:rq1}
\setlength{\tabcolsep}{8pt}
\begin{tabular}{@{}lcc@{}}
\toprule
\thead{ODC type} & \thead{Count} & \thead{Share} \\
\midrule
Algorithm/Method            & \TBD & \TBD \\
Assignment/Initialization   & \TBD & \TBD \\
Checking                    & \TBD & \TBD \\
Timing/Serialization        & \TBD & \TBD \\
Function/Class/Object       & \TBD & \TBD \\
Interface/O-O Messages      & \TBD & \TBD \\
Relationship                & \TBD & \TBD \\
\midrule
\textbf{Total}              & \TBD & \TBD \\
\bottomrule
\end{tabular}
\end{table}

As a descriptive companion, the same prefix pass also yields an Impact
distribution (Section~\ref{sec:impact}): modal category \TBD{}, label entropy
\TBD{}, and bug-report coverage \TBD{}\% (uneven across projects;
see Section~\ref{sec:threats}). This is reported for completeness, not as a claim
of discriminative power.

\subsection{RQ2: Taxonomy Coverage}

RQ2 asks how much of \textit{Defects4J} the seven ODC types actually cover. The
open-set formulation makes this measurable: the \textit{scientific-open} pass may
route a bug to ``Other'' (with mandatory justification), and the escape rate is
the share of bugs that do so. We also run the \textit{scientific-closed} pass on
the same bugs and measure the closed$\leftrightarrow$open shift with Cohen's
$\kappa$, quantifying how often forcing a choice changes the assigned type.
Table~\ref{tab:rq2} reports these coverage metrics.

\begin{table}[H]
\centering
\caption{RQ2: taxonomy coverage (\textit{scientific-closed} vs.\
         \textit{scientific-open}).}
\label{tab:rq2}
\setlength{\tabcolsep}{8pt}
\begin{tabular}{@{}lc@{}}
\toprule
\thead{Metric} & \thead{Value} \\
\midrule
Escape rate (share routed to ``Other'')   & \TBD \\
Coverage rate (share within the seven types) & \TBD \\
Closed$\leftrightarrow$open shift-$\kappa$   & \TBD \\
Bugs whose type shifted                      & \TBD \\
\bottomrule
\end{tabular}
\end{table}

\subsection{RQ3: Classification Accuracy}

RQ3 evaluates accuracy with the four-tier framework of
Section~\ref{sec:eval}, comparing each pre-fix classification against its
\textit{scientific-open} post-fix reference. Table~\ref{tab:rq3} reports the four
tiers. Strict match is the exact-agreement rate; top-2 and family match relax
agreement to the alternative list and the ODC family respectively; and
Cohen's~$\kappa$ gives the chance-corrected agreement between the two evidence
modes, interpreted on the Landis--Koch scale \citep{landis1977measurement}.

\begin{table}[H]
\centering
\caption{RQ3: four-tier pre-fix vs.\ post-fix agreement (\textit{scientific-open}).}
\label{tab:rq3}
\setlength{\tabcolsep}{7pt}
\begin{tabular}{@{}llc@{}}
\toprule
\thead{Tier} & \thead{Metric} & \thead{Value} \\
\midrule
1 & Strict type match     & \TBD \\
2 & Top-2 match           & \TBD \\
3 & Family match          & \TBD \\
4 & Cohen's $\kappa$      & \TBD \\
\bottomrule
\end{tabular}
\end{table}

\subsection{RQ4: Contribution of Each Component}

RQ4 isolates what taxonomy grounding and the reasoning loop each contribute by
walking an ablation ladder from the unstructured baseline
(\textit{zero-free}) through the strong static prompt (\textit{few-open}) to the
enforced loop (\textit{scientific-open}), all pre-fix. For each rung we report the
number of distinct labels produced, the label entropy, and how many of the seven
ODC types are exercised; the vocabulary-reduction ratio summarises how sharply a
taxonomy collapses free-form labels. Table~\ref{tab:rq4} reports the ladder.

\begin{table}[H]
\centering
\caption{RQ4: ablation ladder (pre-fix), from unstructured baseline to the
         enforced loop.}
\label{tab:rq4}
\setlength{\tabcolsep}{5pt}
\begin{tabular}{@{}lccc@{}}
\toprule
\thead{Condition} & \thead{Unique labels} & \thead{Entropy} & \thead{ODC coverage} \\
\midrule
zero-free        & \TBD & \TBD & \TBD \\
few-open         & \TBD & \TBD & \TBD \\
scientific-open  & \TBD & \TBD & \TBD \\
\midrule
\multicolumn{3}{@{}l}{Vocabulary-reduction ratio} & \TBD \\
\bottomrule
\end{tabular}
\end{table}

\subsection{RQ5: Effect of Seeing the Fix}

RQ5 quantifies how much the available evidence, rather than the bug itself,
shapes the assigned type, and how reliability varies across projects. Globally it
reuses the pre/post agreement of RQ3; Table~\ref{tab:rq5} adds the per-project
Cohen's~$\kappa$, which requires at least two bugs per project to be defined and
is therefore a primary driver of the per-project minimum in the confirmatory
run's manifest.

\begin{table}[H]
\centering
\caption{RQ5: per-project pre-fix/post-fix reliability (Cohen's~$\kappa$,
         \textit{scientific-open}).}
\label{tab:rq5}
\setlength{\tabcolsep}{10pt}
\begin{tabular}{@{}lc@{}}
\toprule
\thead{Project} & \thead{$\kappa$} \\
\midrule
Chart & \TBD \\ Cli & \TBD \\ Closure & \TBD \\ Codec & \TBD \\
Collections & \TBD \\ Compress & \TBD \\ Csv & \TBD \\ Gson & \TBD \\
JacksonCore & \TBD \\ JacksonDatabind & \TBD \\ JacksonXml & \TBD \\
Jsoup & \TBD \\ JxPath & \TBD \\ Lang & \TBD \\ Math & \TBD \\
Mockito & \TBD \\ Time & \TBD \\
\bottomrule
\end{tabular}
\end{table}

\subsection{Interpreting Divergence}
\label{sec:divergence}

Pre-fix and post-fix classifications are expected to diverge for a minority of
bugs, for reasons intrinsic to semantic classification rather than pipeline
error. \textbf{Evidence asymmetry}: pre-fix evidence shows what went wrong
(symptoms) while post-fix evidence shows what was changed (the fix), and one
underlying defect can produce several symptoms. \textbf{ODC boundary ambiguity}: some bugs
legitimately sit near a type boundary, most often between \textit{Checking}
and \textit{Algorithm/Method}, where a missing guard and a corrected computation
describe the same defect from two angles. \textbf{Multi-fault contamination}:
co-existing faults \citep{callaghan2024mf} can inject symptoms from a different
fault than the one being fixed. The four-tier framework is designed to make this
structure legible: a gap between strict and family match localises disagreement
to boundary-adjacent types rather than to unrelated families. A representative
soft-divergence case study will be populated from the confirmatory run.

\textbf{Impact as a negative control.} Because Impact is opener-defined and
fix-independent (Section~\ref{sec:impact}), its pre-fix/post-fix agreement over
the same paired bugs should reflect instrument noise alone, with no genuine
fix-dependence to explain disagreement. We report it alongside Defect-Type
drift as a noise floor: raw agreement \TBD{} and, where the marginal
distribution is non-degenerate, Cohen's~$\kappa$ \TBD{}. Defect-Type drift
materially above this floor indicates genuine fix-dependence in the assigned
type; drift at or below it would instead suggest the apparent semantic
divergence is largely instrument noise.

%% =============================================================================
